\documentclass[12pt]{amsart} 
\usepackage{mathtools}

\input{./latex-preamble/cm-preamble.tex}
\input{./latex-preamble/use-scr.tex}
\input{./latex-preamble/fr-mod.tex}

%\graphicspath{ {./diagrams/} }

\begin{document}

{Hartshorne} \hfill {9-17-2026}

\bigskip\bigskip

\begin{center}

{\sc Chapter III.7}

\end{center}

\begin{center}

  {Noah Parker} 
    
\end{center}

\bigskip\bigskip

\begin{enumerate}
  \item[II.8.1.] {
      Here we will strengthen the results of the text to include information about the sheaf of differentials at a not necessarily closed point of a scheme $X$.
      \begin{enumerate}
        \item Generalize (8.7) as follows.
          Let $B$ be a local ring containing a field $k$, and assume that the residue field $k(B)=B/\m$ of $B$ is a separably generated extension of $k$.
          Then the exact sequence of (8.4A),
          \[
            0\to \m/\m^2\xrightarrow\delta \Omega_{B/k}\otimes k(B)\to \Omega_{k(B)/k}\to 0
          \] 
          is exact on the left also.
          [\emph{Hint}: In copying the proof of (8.7), first pass to $B/\m^2$, which is a complete local ring, and then use (8.25A) to choose representatives for $B/\m^2$.]
        \item Generalize (8.8) as follows.
          With $B$, $k$ as above, assume furthermore that $k$ is perfect, and that $B$ is a localization of an algebra of finite type over $k$.
          Then show that $B$ is a regular local ring if and only if $\Omega_{B/k}$ is free of rank $=\dim B+\mb{tr.d.}\,k(B)/k$.
        \item Strengthen (8.15) as follows.
          Let $X$ be an irreducible scheme of finite type over a perfect field $k$, and let $\dim X=n$.
          For any point $x\in X$, not necessarily closed, show that the local ring $\O_{X,x}$ is a regular local ring if and only if the stalk $(\Omega_{X/k})_x$ of the sheaf of differentials at $x$ is free of rank $n$.
        \item Strengthen (8.16) as follows.
          If $X$ is a variety over and algebraically closed field $k$, then $U=\{x\in X|\O_x\textup{ is a regular local ring}\}$ is an open dense subset of $X$.
      \end{enumerate}
    }
\end{enumerate}

\begin{enumerate}
  \item[II.8.6.] {
      \emph{The Infinitesimal Lifting Prperty}.
    }
\end{enumerate}

\begin{enumerate}
  \item[6.10.] {
      \emph{Duality for a Finite Flat Morphism.}
      \begin{enumerate}
        \item Let $f:X\to Y$ be a finite morphism of noetherian schemes.
          For any quasi-coherent $\O_Y$-module $\G$, $\sHom_Y(f_*\O_Y,\G)$ is a quasi-coherent $f_*\O_Y$-module, and hence corresponds to a quasi-coherent $\O_X$-module, which we call $f^!\G$ (II, Ex. 5.17e).
        \item Show that for any coherent $\F$ on $X$ and any quasi-coherent $\G$ on $Y$, there is a natural isomorphism
          \[
            f_*\sHom_X(\F,f^!\G)\iso \sHom_Y(f_*\F,\G).
          \] 
        \item For each $i\geqs 0$, there is a natural map
          \[
            \phi_i:\Ext^i_X(\F,f^!\G)\to \Ext^i_Y(f_*\F,\G)
          \] 
          [\emph{Hint}: First construct a map
          \[
            \Ext^i_X(\F,f^!\G)\to \Ext^i_Y(f_*\F,f_*f^!\G).
          \] 
          Then compose with a suitable map from $f_*f^!\G\to G$.]
        \item Now assume that $X$ and $Y$ are separated, $\Coh(X)$ has enough locally frees, and assume that $f_*\O_X$ is locally free on $Y$ (this is equivalent to saying $f$ is flat--see Section 9).
        Show that $\phi_i$ is an isomorphism for all $i$, all $\F$ coherent on $X$, and all $\G$ quasi-coherent on $Y$.
        [\emph{Hints}: First do $i=0$.
        Then do $\F=\O_X$, using (Ex. 4.1).
        Then do $\F$ locally free.
        Do the general case by induction on $i$, writing $\F$ as a quotient of a locally free sheaf.]
      \end{enumerate}
    }
\end{enumerate}

\begin{proof}
\end{proof}

\begin{enumerate}
  \item[7.2.] {
      Let $f:X\to Y$ be a finite morphism of projective schemes over a field $k$, and let $\omega_Y^\circ$ be a dualizing sheaf for $Y$.
      \begin{enumerate}
        \item Show that $f^!\omega_Y^\circ$ is a dualizing sheaf for $X$, where $f^!$ is defined as in (Ex. 6.10).
        \item If $X$ and $Y$ are both nonsingular, and $k$ is algebraically closed, conclude that there is a natural trace map $t:f_*\omega_X\to \omega_Y$.
      \end{enumerate}
    }
\end{enumerate}

\begin{proof}
  (a) Fix a coherent sheaf $\F\in \Mod(X)$, and a trace map $t_Y:H^n(Y,\omega_Y^\circ)\to k$. 
  (Ex. 6.10c) Gives us a map
  \begin{align*}
    \phi_0:\Hom_X(\F,f^!\omega_Y^\circ)\iso \Hom_Y(f_*\F,\omega_Y^\circ),
  \end{align*}
  which is
  $f_*\F$ is coherent by (Ex. II.5.5), so there is a map
  \begin{align*}
    \Hom_Y(f_*\F,\omega_Y^\circ)\times H^n(Y,f_*\F)\to H^n(Y,\omega_Y^\circ)
  \end{align*}
  which, composed with $t_Y:H^n(Y,\omega_Y^\circ)\to k$, induces a perfect pairing.

  (b) Part (a), (Ex. 6.10b), and give us isomorphisms
  \begin{align*}
    f_*\omega_X&\cong f_*\sHom_X(\O_X,\omega_X) \\
               &\cong f_*\sHom_X(\O_X,\omega_Y)
  \end{align*}
\end{proof}

\begin{enumerate}
  \item[7.3.] {
      Let $X=\P_k^n$.
      Show that $H^q(X,\Omega_X^p)=0$ for $p\neq q$, $k$ for $p=q, 0\leqs p,q\leqs n$.
    }
\end{enumerate}

\begin{proof}
\end{proof}

\begin{enumerate}
  \item[7.4.] {
      \emph{The Cohomology Class of a Subvariety.}
      Let $X$ be a nonsingular projective variety of dimension $n$ over an algebraically closed field $k$.
      Let $Y$ be a nonsingular subvariety of codimension $p$ (hence dimension $n-p$).
      From the natural map $\Omega_X\otimes \O_Y\to \Omega_Y$ (II, 8.12) we deduce a map $\Omega_X^{n-p}\to \Omega_Y^{n-p}$.
      This induces a map on cohomology $H^{n-p}(X,\Omega_X^{n-p})\to H^{n-p}(Y,\Omega_Y^{-n-p})$.
      Now $\Omega_Y^{n-p}=\omega_Y$ is a dualizing sheaf for $Y$, so we have the trace map $t_Y:H^{n-p}(Y,\Omega_Y^{n-p})\to k$.
      Composing, we obtain a linear map $H^{n-p}(X,\Omega_X^{n-p})\to k$.
      By (7.13) this corresponds to an element $\eta(Y)\in H^p(X,\Omega_X^p)$, which we call the \emph{cohomology clas} of $Y$.
      \begin{enumerate}
        \item If $P\in X$ is a closed point, show that $t_X(\eta(P))=1$, where $\eta(P)\in H^n(X,\Omega^n)$, and $t_X$ is the trace map.
        \item If $X=\PP^n$, identify $H^p(X,\Omega^p)$ with $k$ by (Ex. 7.3), and show that $\eta(Y)=(\deg Y)\cdot 1$, where $\deg Y$ is its \emph{degree} as a projective variety.
          [\emph{Hint}: Cut with hyperplanes $H\subset X$, and use Bertini's Theorem (II, 8.18) to reduce to the  case $Y$ is a finite set of points.]
        \item For any scheme $X$ of finite type over $k$, we define a homomorphism of sheaves of abelian groups $d\log:\O_X^\times\to \Omega_X$ by $d\log(f)=f^{-1}df$.
          Here $\O^\times$ is a group under multiplication, and $\Omega_X$ is a group under addition.
          This induces a map on cohomology $\Pic X=H^1(X,\O_X^\times)\to H^1(X,\Omega_X)$, which we denote by $c$-see (Ex. 4.5).
        \item Returning to the hypotheses above, suppose $p=1$.
          Show that $\eta(Y)=c(\L(Y))$, where $\L(Y)$ is the invertible sheaf corresponding to the divisor $Y$.
      \end{enumerate}
    }
\end{enumerate}

\begin{proof}
  (a)
\end{proof}

\end{document}
